\documentclass[10pt,a4paper]{amsart}
\usepackage[margin=19mm]{geometry}
\usepackage{amsmath,amssymb,mathtools}
\usepackage[scr=rsfs,cal=euler]{mathalfa}
\usepackage{xfrac}
\usepackage[unicode,hidelinks,bookmarks=false]{hyperref}
\newcommand{\ie}{i.e.\ }
\newcommand{\cf}{cf.\ }
\newcommand{\resp}{respectively\ }
\newcommand{\Subsection}{\subsection}
\providecommand{\nobreakdash}{\nobreak\hbox{-}}
\setlength{\emergencystretch}{2em}
\multlinegap0pt
\usepackage{float}
\newtheorem{theorem}{Theorem}
\newcommand{\ceil}[1]{\left\lceil #1\right\rceil}
\newcommand{\disc}{\operatorname{disc}}
\newcommand{\floor}[1]{\left\lfloor #1\right\rfloor}
\title{The Asayama-Matsumoto conjecture and a refined discrepancy bound}
\author{Alessio Basti, Tommaso Cremaschi}
\date{Source: arXiv:2608.21585v1 (CC BY 4.0)}
\begin{document}
\maketitle

\noindent\emph{Source and licence.} The Asayama-Matsumoto conjecture and a refined discrepancy bound. Version arXiv:2608.21585v1 (CC BY 4.0), distributed by arXiv under the Creative Commons Attribution 4.0 International licence, http://creativecommons.org/licenses/by/4.0/, as recorded in the arXiv licence metadata for this exact version and shown on the abstract page https://arxiv.org/abs/2608.21585v1. Full licence notice: \texttt{licenses/arxiv-2608.21585-CC-BY-4.0.txt}.\par\smallskip

\noindent\textit{Editorial scope.} Counted unit: the article's theorem thm:refined (source0.tex lines 196--203). The article's complete original proof of it is delivered with it (source0.tex lines 205--284, one \texttt{proof} environment of 80 lines). No mathematics is written, repaired or re-derived: the statement and the proof are the article's own lines, in the article's own notation, and no retained line is substituted or re-flowed.  Retained uncounted as the source-local context the proof needs: lines 107--114.  Nothing was omitted from any retained span, so the package carries no omission record.  The retained lines cite 2 works, whose entries are transcribed verbatim from the article's own bibliography source in the e-print and delivered with short editorial labels recorded in the item's reference mapping.  No retained line contains a figure environment, an image-inclusion command or drawing code, so no image is delivered and the package declares no asset dependency.  Editorial formatting only, in addition: an AMS \texttt{amsart} preamble that restates the article's own declarations and macros used by the retained lines, namely the environments \texttt{theorem} on separate counters, together with the standard packages those lines need.  The retained environments print in this document as Theorem 1 in order of appearance, whereas the article prints the same items with the numbering of its own full document.  The complete native source of the article as distributed in the arXiv e-print for this exact version is bundled unchanged as \texttt{sources/208259/source0.tex}.
\begin{theorem}[Asayama--Matsumoto conjecture]\label{thm:universal}
Every plane triangulation $T$ on $n\ge3$ vertices satisfies
\begin{equation}\label{eq:universal}
  \disc(T)\le n-2\ceil{\frac n3}\le\floor{\frac n3}.
\end{equation}
Moreover, a coloring whose discrepancy satisfies the first bound can be found in
$O(n\log n)$ time.
\end{theorem}

\begin{theorem}[Refined bound]\label{thm:refined}
Let $T$ be a plane triangulation on $n\ge6$ vertices.  If
$n\not\equiv5\pmod6$, then
\begin{equation}\label{eq:refined}
  \disc(T)\le n-2\ceil{\frac{n+2}{3}}.
\end{equation}
Moreover, a coloring satisfying this estimate can be found in $O(n\log n)$ time.
\end{theorem}

\begin{proof}
Take a proper $4$-coloring supplied by
\cite[Corollary~17]{KawarabayashiEtAl2026}, and order its color classes
$V_1,V_2,V_3,V_4$ so that, writing
\[
  a=|V_1|\ge b=|V_2|\ge c=|V_3|\ge d=|V_4|,
\]
we have
\begin{equation}\label{eq:balanced-four-refined}
  a\le\ceil{\frac{n-2}{2}}<\frac n2.
\end{equation}
Set
\[
  R=V_1\cup V_4,
  \qquad
  B=V_2\cup V_3,
  \qquad
  r=a+d,
  \qquad
  s=b+c=n-r.
\]
As in the proof of Theorem~\ref{thm:universal}, this merge is
polychromatic.  Moreover, $r\ge b,c$ gives $r\ge n/3$, while
$d\le b,c$ and~\eqref{eq:balanced-four-refined} give
\[
  r=a+d\le\frac{n+2a}{3}<\frac{2n}{3}.
\]
Thus
\begin{equation}\label{eq:baseline-both}
  r,s\ge\ceil{\frac n3}.
\end{equation}
We show that the smaller class in fact has size at least
$\ceil{(n+2)/3}$.

Suppose first that $n=3q$.  By~\eqref{eq:baseline-both}, $r\ge q$,
and $r<2q$ implies $s\ge q+1$.  If $r\ge q+1$, we are done.  If
$r=q$, then $b,c\le r$ and $b+c=2q$, so $b=c=q$.  Since $a\ge b$
and $a+d=q$, we obtain $a=q$ and $d=0$.  Every face therefore contains
one vertex from each of $V_1,V_2,V_3$.  Color $V_2$ red, $V_3$ blue,
one vertex of $V_1$ red, and the remaining vertices of $V_1$ blue.  The
coloring is polychromatic and its two classes have sizes $q+1$ and
$2q-1$, both at least $q+1$ because $q\ge2$.  This direct split is a special case of the general result of Asayama and
Matsumoto that every properly $3$-colorable triangulation has a balanced
polychromatic $2$-coloring~\cite{AsayamaMatsumoto2022}.

If $n=3q+1$, then
\[
  \ceil{\frac n3}=q+1=\ceil{\frac{n+2}{3}},
\]
so Theorem~\ref{thm:universal} already gives the claim.

It remains to consider $n=3q+2$.  The assumption
$n\not\equiv5\pmod6$ says that $q$ is even, and $q\ge2$ because
$n\ge6$.  By~\eqref{eq:baseline-both}, $r,s\ge q+1$.

If $r=q+1$, then $b+c=2q+1$.  Since $b,c\le r$ and $b\ge c$, we get
$b=q+1$ and $c=q$.  The inequalities $a\ge b$ and $a+d=q+1$ then give
$a=q+1$ and $d=0$.  Again every face contains one vertex from each of
$V_1,V_2,V_3$.  Color $V_2$ red, $V_3$ blue, one vertex of $V_1$ red,
and the remaining vertices of $V_1$ blue.  The two classes have sizes
$q+2$ and $2q$, both at least $q+2$.

Finally, suppose that $s=q+1$, so $r=2q+1$.  Since $q$ is even,
\eqref{eq:balanced-four-refined} gives $a\le3q/2$, and hence
\[
  d=r-a\ge\frac q2+1.
\]
As $b,c\ge d$, it follows that
\[
  s=b+c\ge2d\ge q+2,
\]
a contradiction.  Thus neither merged class has size $q+1$, and both
have size at least
$q+2=\ceil{(n+2)/3}$.  This proves~\eqref{eq:refined}.

The construction consists of the $O(n\log n)$ balanced four-coloring
algorithm, class-size comparisons, the merge above, and, in the two
three-color cases, the choice of a single vertex.  Its total running time
is therefore $O(n\log n)$.
\end{proof}

\begin{thebibliography}{99}
\bibitem[Asayam]{AsayamaMatsumoto2022} Y.~Asayama and N.~Matsumoto, \emph{Balanced polychromatic $2$-coloring of triangulations}, Graphs Combin. \textbf{38} (2022), Article~16, 12~pp., \href{https://doi.org/10.1007/s00373-021-02420-8} {doi:10.1007/s00373-021-02420-8}.
\bibitem[Kawara]{KawarabayashiEtAl2026} K.-i.~Kawarabayashi, H.~Yoneda, and M.~Yoneda, \emph{The balanced four-color theorem}, arXiv:2607.13025v1 [cs.DS], July 14, 2026, \href{https://doi.org/10.48550/arXiv.2607.13025} {doi:10.48550/arXiv.2607.13025}.
\end{thebibliography}
\end{document}
